\VerdierPhysicalPageAnchor{0039}%
\VerdierSemanticAnchor{ch1-1-2-8}
\noindent\textbf{Définition 1.2.8.}\quad
Soient $(\mathcal C_e,G^\circ,T,c)$ et
$(\mathcal C'_e,G^\circ,T',c')$ deux $G^\circ$-catégories et
$(F_e,m)$ et $(F'_e,m')\colon\mathcal C_e\to\mathcal C'_e$ deux
$G^\circ$-foncteurs. On dit qu'un morphisme de foncteurs
$u_e\colon F_e\to F'_e$ est \emph{compatible avec les opérations de
$G^\circ$} si pour tout $g\in G$ le diagramme (1.2.7.2) est commutatif.

\VerdierSemanticAnchor{ch1-1-2-additive}
Enfin, nous dirons qu'une $G^\circ$-catégorie est \emph{additive} si la
catégorie sous-jacente est additive. D'après (1.1.3), la
$G^\circ$-catégorie associée à une catégorie graduée additive est une
$G^\circ$-catégorie additive.

\VerdierSemanticAnchor{ch1-1-3}
\subsection{Équivalence de deux langages}

\VerdierSemanticAnchor{ch1-1-3-1}
\noindent\textbf{1.3.1.}\quad
On se fixe un univers $U$. Désignons par $\operatorname{Grad}(G)$ la
$U$-catégorie dont les objets sont les catégories graduées de type $G$,
appartenant à $U$, et dont les morphismes sont les foncteurs gradués.
Désignons de même par $G^\circ\text{-}\operatorname{Cat}$ la
$U$-catégorie dont les objets sont les $G^\circ$-catégories, appartenant
à $U$, et dont les morphismes sont les $G^\circ$-foncteurs. On déduit
immédiatement de (1.2.1), (1.2.3) et (1.2.5) que le choix, pour toute
catégorie graduée de type $G$, d'un clivage normalisé, définit un
foncteur :

\VerdierSemanticAnchor{ch1-1-3-1-eq1}
\begin{equation}
  \Psi\colon \operatorname{Grad}(G)
  \longrightarrow G^\circ\text{-}\operatorname{Cat}.
  \tag{1.3.1.1}
\end{equation}

Le résultat de [18] s'interprète comme suit :

\VerdierSemanticAnchor{ch1-1-3-2}
\noindent\textbf{Proposition 1.3.2.}\quad
\emph{Le foncteur $\Psi$ est une équivalence de catégories.}

\VerdierSemanticAnchor{ch1-1-3-2-quasi-inverse}
On trouvera dans [18] la description d'un foncteur quasi-inverse de
$\Psi$. Nous nous bornerons ici à décrire rapidement la catégorie
graduée $\mathcal C[G]$ associée à une $G^\circ$-catégorie
$(\mathcal C,G^\circ,T,c)$. Tout d'abord, l'ensemble des objets de
$\mathcal C[G]$ est l'ensemble des objets de $\mathcal C$. Soient $X$ et
$Y$ deux objets de $\mathcal C$. Posons alors :

\VerdierSemanticAnchor{ch1-1-3-2-eq1}
\begin{equation}
  \operatorname{Hom}^{g}(X,Y)
  =\operatorname{Hom}_{\mathcal C}(X,T(g)Y),
  \qquad g\in G,
  \tag{1.3.2.1}
\end{equation}

\VerdierSemanticAnchor{ch1-1-3-2-eq2}
\begin{equation}
  \operatorname{Hom}_{\mathcal C[G]}(X,Y)
  =\coprod_{g\in G}\operatorname{Hom}^{g}(X,Y).
  \tag{1.3.2.2}
\end{equation}

\VerdierSemanticAnchor{ch1-1-3-2-composition}
Décrivons alors la composition des morphismes de $\mathcal C[G]$.
Soient $X$, $Y$ et $Z$ trois objets de $\mathcal C$ et
$u\in\operatorname{Hom}^{g}(X,Y)$,
$v\in\operatorname{Hom}^{g'}(Y,Z)$, $g,g'\in G$. Le foncteur $T(g)$
définit une application :

\VerdierSemanticAnchor{ch1-1-3-2-eq3}
\begin{equation}
  \operatorname{Hom}^{g'}(Y,Z)
  \longrightarrow
  \operatorname{Hom}_{\mathcal C}
    \bigl(T(g)Y,T(g)T(g')Z\bigr)
  \tag{1.3.2.3}
\end{equation}
